\documentclass[11pt]{article}
\usepackage[letterpaper,margin=1in]{geometry}
\usepackage{times}
\usepackage{amsmath}
\pagestyle{empty}
\begin{document}

\begin{center}
{\Large\bfseries Device-normalised geometry for muscle architecture\\[4pt]
measurement in B-mode ultrasound}
\end{center}
\vspace{1.5em}

\noindent\textbf{Author}
\vspace{0.4em}

\noindent Colton Downing\\
University of Southern California, Los Angeles, CA, USA\\
ORCID: 0009-0009-8874-1914
\vspace{1.2em}

\noindent\textbf{Corresponding author}
\vspace{0.4em}

\noindent Colton Downing\\
University of Southern California\\
University Park Campus, Los Angeles, CA 90089, United States\\
cjdownin@usc.edu
\vspace{1.2em}

\noindent\textbf{Declarations}
\vspace{0.4em}

\noindent Conflicts of interest: none. Funding: none. The author declares no competing
financial or personal interests.
\vspace{1.5em}

\hrule
\vspace{1.5em}

\noindent\textbf{\large Abstract}
\vspace{0.8em}

Tools that measure muscle architecture from B-mode ultrasound assume a known
setup: one scanner, one export convention, an operator-supplied scale. This
paper handles the case where none of that exists. The input is an unlabelled
screenshot from one of six ultrasound systems, and pennation angle (PA),
fascicle length (FL) and muscle thickness (MT) have to come out of the pixels
alone.

Two corrections do most of the work, and both follow from first principles
rather than from fitting. The first concerns pixel shape: two of the six
layouts export a horizontally resampled image, so the pixels are not square and
any angle read in pixel coordinates is wrong by a fixed factor. The published
transducer aperture gives that factor directly. The second concerns fascicle
length. The usual estimator extends a fascicle fragment until it meets both
aponeuroses, which is exact only when they run parallel. When they converge with
half-angle $h$, that chord exceeds $\mathrm{MT}/\sin\gamma$ by a factor
$1/(1-\cot^{2}\gamma\tan^{2}h)$. Expert raters do not reproduce that inflation.
Measuring the angle to the bisector of the two aponeuroses removes the
dependence on convergence and lets a single calibration constant carry between
muscles.

On a 35-image benchmark measured independently by seven expert raters, the
method scores 0.196 under the UMUD error formula. DL\_Track\_US scores 0.331,
SMA scores 0.409, and an individual expert rater scores 0.264 by leave-one-out.
On a separate 167-frame dynamic sequence of a different muscle, fascicle-length
error is 1.82~mm using a constant fitted only on the static benchmark.

\vspace{1.2em}

\noindent\textbf{Corresponding author}
\vspace{0.4em}

\noindent Colton Downing\\
University of Southern California\\
University Park Campus, Los Angeles, CA 90089, United States\\
cjdownin@usc.edu
\vspace{1.2em}

\noindent\textbf{Declarations}
\vspace{0.4em}

\noindent Conflicts of interest: none. Funding: none. The author declares no competing
financial or personal interests.
\vspace{1.5em}

\hrule
\vspace{1.5em}

\noindent\textbf{\large Abstract}
\vspace{0.8em}

Tools that measure muscle architecture from B-mode ultrasound assume a known
setup: one scanner, one export convention, and an operator-supplied scale. This
paper handles the case where none of that exists. The input is an unlabelled
screenshot from one of six ultrasound systems, carrying no metadata and no
stated scale, and pennation angle (PA), fascicle length (FL) and muscle
thickness (MT) have to come out of the pixels alone.

Two corrections do most of the work, and both follow from first principles
rather than from fitting. The first concerns pixel shape: two of the six layouts
export a horizontally resampled image, so the pixels are not square and any
angle read in pixel coordinates is wrong by a fixed factor. The published
transducer aperture gives that factor directly (Juniper 12L3,
$R \approx 1.12$; Lumify L12-4, $R \approx 0.84$). The second concerns fascicle
length. The usual estimator extends a fascicle fragment until it meets both
aponeuroses, which is exact only when they run parallel. When they converge with
half-angle $h$, that chord exceeds $\mathrm{MT}/\sin\gamma$ by a factor
$1/(1-\cot^{2}\gamma\tan^{2}h)$. Expert raters do not reproduce that inflation.
Measuring the angle to the bisector of the two aponeuroses removes the
dependence on convergence and lets a single calibration constant carry between
muscles.

On a 35-image benchmark measured independently by seven expert raters, the
method scores 0.196 under the UMUD error formula. DL\_Track\_US scores 0.331 on
the same images, SMA scores 0.409, and an individual expert rater scores 0.264
by leave-one-out. On a separate 167-frame dynamic sequence of a different
muscle, fascicle-length error is 1.82~mm using a constant fitted only on the
static benchmark.

\vspace{1.2em}

\noindent\textbf{Corresponding author}
\vspace{0.4em}

\noindent Colton Downing\\
University of Southern California\\
University Park Campus, Los Angeles, CA 90089, United States\\
cjdownin@usc.edu
\vspace{1.2em}

\noindent\textbf{Declarations}
\vspace{0.4em}

\noindent Conflicts of interest: none. Funding: none. The author declares no competing
financial or personal interests.
\vspace{1.5em}

\hrule
\vspace{1.5em}

\noindent\textbf{\large Abstract}
\vspace{0.8em}

\noindent Tools that measure muscle architecture from B-mode ultrasound usually assume a
known setup: one scanner, one export convention, and a scale the operator types
in. This paper handles the case where none of that exists. The input is an
unlabelled screenshot from one of six different ultrasound systems, carrying no
metadata and no stated scale, and pennation angle (PA), fascicle length (FL) and
muscle thickness (MT) have to come out of the pixels alone.

\vspace{0.6em}
\noindent Two corrections do most of the work, and both follow from first principles
rather than from fitting. The first concerns pixel shape. Two of the six layouts
export a horizontally resampled image, so the pixels are not square, and any
angle read in pixel coordinates on those devices is wrong by a fixed factor. The
published transducer aperture gives that factor directly (Juniper 12L3,
$R \approx 1.12$; Lumify L12-4, $R \approx 0.84$), so the pipeline resamples to
square physical pixels before measuring anything. The second concerns fascicle
length. The usual estimator extends a detected fascicle fragment until it meets
both aponeuroses, which is exact only when the aponeuroses run parallel. When
they converge with half-angle $h$, that chord exceeds $\mathrm{MT}/\sin\gamma$
by a factor $1/(1-\cot^{2}\gamma\tan^{2}h)$. Expert raters do not reproduce that
inflation. Measuring the angle to the bisector of the two aponeuroses instead
removes the dependence on convergence, which cuts the spread in bias across two
different muscles from six percentage points to one and lets a single
calibration constant carry between them.

\vspace{0.6em}
\noindent On a 35-image benchmark measured independently by seven expert raters, the
method scores 0.196 under the UMUD error formula against the seven-rater
consensus. DL\_Track\_US scores 0.331 on the same images, SMA scores 0.409, and
an individual expert rater scores 0.264 by leave-one-out. On a separate
167-frame dynamic sequence of a different muscle, fascicle-length error is
1.82~mm using a constant fitted only on the static benchmark.

\vspace{1.2em}
\noindent\textbf{Keywords:} musculoskeletal ultrasound; muscle architecture; fascicle
length; pennation angle; image calibration

\end{document}
